\section{CPU Architecture Overview}

\subsection{The basic CPU loop}

Before looking at individual modules, it helps to understand the life of a typical instruction.

For example:

\begin{lstlisting}[language={}]
add x3, x1, x2
\end{lstlisting}

Conceptually, the CPU does this:

\begin{lstlisting}[language={}]
1. Use PC to fetch a 32-bit instruction word.
2. Decode that word and discover it is an R-type ADD.
3. Read source registers x1 and x2.
4. Send both values to the ALU.
5. ALU computes rs1 + rs2.
6. Write the result into destination register x3.
7. Advance PC to PC + 4.
\end{lstlisting}

A load instruction adds a memory read. A store adds a memory write. A branch changes the next PC depending on a comparison. A trap redirects the PC to a trap handler.

\subsection{Main conceptual block diagram}

\begin{lstlisting}[language={},caption={Mermaid diagram source}]
flowchart LR
    PC[Program Counter] --> IF[Instruction Fetch]
    IF --> DEC[Instruction Decoder]
    DEC --> REG[Register File]
    DEC --> IMM[Immediate Generator]
    REG --> ALU[ALU]
    IMM --> ALU
    ALU --> BR[Branch / Jump Decision]
    ALU --> LSU[Load / Store Address]
    LSU --> MEM[Data Memory / Bus / MMU]
    MEM --> WB[Write Back]
    ALU --> WB
    WB --> REG
    BR --> PC
    DEC --> CTRL[Control Signals]
    CTRL --> ALU
    CTRL --> LSU
    CTRL --> WB
\end{lstlisting}

\subsection{Single-cycle vs multi-cycle vs pipelined}

The repository contains all three styles.

\subsubsection{Single-cycle}

In a single-cycle CPU, one instruction completes in one clock period. The clock period must be long enough for the worst instruction path: instruction memory, decode, register read, ALU, memory, write-back mux, and setup to the register file.

This is easy to understand but inefficient for FPGA block RAM and high clock frequency. In this project, \texttt{rtl/cpu.v} is the cleanest self-contained single-cycle CPU, while \texttt{rtl/cpu\_core.v} exposes the data-memory side as an external bus for SoC use.

\subsubsection{Multi-cycle}

In a multi-cycle CPU, the same instruction is split across several clock cycles. The CPU uses a small finite-state machine: fetch the instruction, execute it, optionally access memory, then return to fetch. This is slower in instructions per clock but easier to map to registered memories.

\texttt{rtl/cpu\_mc.v} is this version. The comments describe states such as \texttt{FETCH}, \texttt{EXEC}, and \texttt{MEM}: ALU/branch/store instructions take fewer states than loads, because loads need an extra memory-data cycle.

\subsubsection{Pipelined}

In a pipelined CPU, multiple instructions are in flight at once. A classic five-stage pipeline uses:

\begin{lstlisting}[language={}]
IF  -> instruction fetch
ID  -> decode and register read
EX  -> execute / branch decision / address calculation
MEM -> data memory access
WB  -> register write-back
\end{lstlisting}

\texttt{rtl/cpu\_pipe.v} implements this style. It improves throughput, but introduces hazards: one instruction may need a value that an older instruction has not written back yet, or the CPU may fetch down the wrong path after a branch.

\bigskip
